\subNsection{1}{Java-Programmierumgebung}{Markus Nöth}

Die Java-Programmierumgebung ist Bestandteil des Moduls \glqq Programmieren~1\grqq{} und ermöglicht das Schreiben, Kompilieren und Ausführen von Java-Programmen direkt im Browser ohne Server-Backend.

\subNsection{2}{Entscheidungen in der Architektur}{Markus Nöth}

Die Umgebung folgt einer mehrschichtigen Architektur aus Präsentation, Orchestrierung und isolierter Ausführungsumgebung:
\begin{enumerate}
  \item \textbf{Präsentationsschicht:} Der \verb|JavaExecutionProvider| stellt den vom Hook verwalteten Zustand als \verb|JavaExecutionSnapshot| bereit (Status, Konsolenausgaben, Diagnosen und Artefakte). Im Monaco-Editor~\cite{monacoeditor} (\verb|JavaEditorPanel|) begrenzt die Anwendung den Quelltext über \verb|javaEditorLimits.ts| auf 500.000~UTF-16-Codeeinheiten und 10.000~Zeilen und visualisiert Fehlermeldungen direkt im Code.
  \item \textbf{Orchestrierungsschicht:} Der Hook \verb|useJavaExecution| steuert den Lebenszyklus des Web Workers (\verb|javaRunner.worker.ts|) über asynchrone Nachrichten mit eindeutigen Request-\glspl{id}. Eine konfigurierbare Laufzeitüberwachung schützt vor Endlosschleifen: Überschreitet der Auftrag einschließlich Laden, Kompilieren und Ausführen das Zeitlimit, beendet der Hauptthread den Worker mittels \texttt{worker.\allowbreak{}terminate()} und startet eine neue Instanz. Die Programmausführung blockiert damit nicht den Hauptthread.
  \item \textbf{Compiler- und \gls{wasm}-Laufzeit:} Ein untergeordneter Compiler-Worker lädt das vorkompilierte \gls{wasm}-Compilermodul (\texttt{teavm-javac}) samt Klassenbibliotheken als statische Assets nach. Dieses bündelt das OpenJDK-\verb|javac|-Frontend mit dem TeaVM-Backend~\cite{teavm}: Nach der Syntax- und Typprüfung durch \verb|javac| übersetzt TeaVM~\cite{teavm} den erzeugten Bytecode in \gls{wasm}, das anschließend im Runner-Worker ausgeführt wird.
\end{enumerate}



Abbildung~\ref{fig:java_worker_ablauf} zeigt den Ablauf:
\begin{enumerate}
  \item \textbf{Quelltext-Vorverarbeitung:} Echte interaktive Konsoleneingaben während der Programmausführung würden eine zusätzliche Abstimmung zwischen Benutzeroberfläche und Wasm-Laufzeit erfordern. Aus diesem Grund setzt das Modul auf einen leichtgewichtigen Vorprozessor: \verb|inlineScannerInput()| ersetzt erkannte Scanner-Aufrufe vor der Kompilierung per regulärem Ausdruck durch feste Werte aus einem vorbereiteten \verb|ScannerInputCursor|. Unterstützt werden unter anderem ganze Zahlen (\texttt{nextInt()}), Dezimalzahlen (\texttt{nextDouble()}), Wahrheitswerte (\texttt{nextBoolean()}) sowie einzelne Wörter und Textzeilen (\texttt{next()}, \texttt{nextLine()}). Bei gültigen, vorab bereitgestellten Eingaben lassen sich damit einfache Programme zur Eingabe, Berechnung und Ausgabe umsetzen, die die grundlegenden Arbeitsschritte für den Einstieg abdecken. Die Ersetzung folgt jedoch der Reihenfolge im Quelltext: Wiederholte Eingaben in Schleifen oder nur bedingt ausgeführte Eingabeaufrufe werden nicht entsprechend dem tatsächlichen Kontrollfluss ausgewertet.
  \item \textbf{Übersetzung und Ausführung:} Der vorverarbeitete Quelltext wird im Web Worker durch \verb|javac| kompiliert und via TeaVM~\cite{teavm} in ein \gls{wasm}-Modul transformiert, das im Runner-Worker instanziiert wird.
  \item \textbf{Pufferung des Ausgabestroms:} Über das Foreign Function Interface (\gls{ffi}) von WebAssembly registriert \verb|installTeaVmConsole()| die Host-Funktionen \texttt{putcharStdout} und \texttt{putcharStderr} in der Import-Tabelle des Wasm-Moduls. Standard- und Fehlerausgaben von \texttt{System.out} bzw. \texttt{System.err} werden so zeichenweise über die FFI-Grenze an JavaScript übergeben. \texttt{create-Buffered-Worker-Output} bündelt die Zeichen je Ausgabestrom mit einer konfigurierten Verzögerung von 50ms. Der Puffer wird gesendet, sobald der Timer ausgeführt werden kann, bei einem weiteren Zeichen nach Ablauf der Verzögerung oder beim Abschluss der Ausführung. Dies reduziert die Zahl einzelner \texttt{postMessage}-Nachrichten und damit auch die Zahl der Rerenders.
\end{enumerate}


\begin{figure}[H]
  \centering
  \begin{tikzpicture}[
    font=\footnotesize,
    box/.style={draw, rounded corners=2pt, align=center, inner sep=5pt, text width=3.8cm, minimum height=1cm},
    flow/.style={-{Stealth[length=2mm]}, semithick},
    lbl/.style={font=\scriptsize, fill=white, inner sep=2pt}
  ]
    \node[box, fill=black!5] (ui) at (0,0) {Hauptthread\\Editor und Ausführungs-Hook};
    \node[box, fill=black!12] (runner) at (5,0) {Runner-Worker\\Scanner-Vorverarbeitung};
    \node[box, fill=black!5] (compiler) at (10,0) {Compiler-Worker\\javac $\rightarrow$ Bytecode\\TeaVM $\rightarrow$ Wasm};
    \node[box, fill=black!12] (wasm) at (5,-2) {Runner-Worker\\Wasm-Ausführung und Ausgabepuffer};
    \draw[flow] (ui) -- (runner);
    \draw[flow] (runner) -- (compiler);
    \draw[flow] (compiler.south) |- (wasm.east) node[lbl, near end, above] {Wasm};
    \draw[flow] (wasm.west) -| (ui.south) node[lbl, near start, above] {Ausgaben, Status};
  \end{tikzpicture}
  \caption[Java-Ausführung mit zwei Workern]{Java-Ausführung mit getrenntem Runner- und Compiler-Worker und TeaVM~\cite{teavm} als Wasm-Backend.}
  \label{fig:java_worker_ablauf}
\end{figure}

\subNsection{2}{Entscheidungen im UI}{Markus Nöth}

Die Arbeitsoberfläche orientiert sich an vertrauten \glspl{ide} wie Visual Studio Code:
\begin{itemize}
  \item \textbf{Monaco-Editor:} Der Editor bietet Code-Faltung, Volltextsuche (\texttt{Strg+F}) und automatische Einrückung. Compilerfehler mit Positionsangabe, beispielsweise fehlende Semikolons oder inkompatible Typen, erscheinen als rote Marker. Entfernte Scanner-Zeilen werden dabei nicht auf den ursprünglichen Quelltext zurückgerechnet.
  \item \textbf{Eingabebereich für \texttt{System.in}:} Über ein separates Texteingabefeld kann der Benutzer Konsoleneingaben vorab definieren, die der Vorprozessor beim Programmstart verarbeitet.
  \item \textbf{Konsole und Problemliste (\texttt{JavaConsolePanel}):} Standardausgabe, Fehlerausgabe und Diagnosen erscheinen in getrennten Abschnitten. Die Diagnoseliste nennt die Meldung sowie, falls vorhanden, Datei und Position; ein Sprung in den Editor ist dort nicht implementiert.
  \item \textbf{Steuerleiste:} Start und Abbruch liegen beim Simulator; Zurücksetzen und Fokusmodus bietet der gemeinsame Simulationsrahmen.
\end{itemize}

\subNsection{2}{Andere Entscheidungen}{Markus Nöth}

Das Modul enthält vier interaktive Beispielprogramme und 13~Übungsaufgaben (u.\,a. Konsolenausgaben, Schleifen, Variablenberechnungen).
Die Aufgabenprüfung kompiliert den eingegebenen Code und vergleicht seine normalisierte Standardausgabe mit den erwarteten Ausgaben definierter Testfälle. Dies belegt keine allgemeine Korrektheit. Bei Aufgaben mit vorgegebenen Variablenbelegungen ersetzt ein regulärer Ausdruck (\verb|sourceTransform|) die Initialisierungswerte im Code durch alternative Testdaten, um weitere Eingabefälle zu prüfen. Ändert der Lernende jedoch die Bezeichner oder den vom Muster erwarteten Aufbau dieser Vorgabezeilen, greift der Ersetzungsausdruck nicht mehr. In diesem Fall kann die Validierung fehlschlagen, obwohl der geschriebene Algorithmus fachlich korrekt ist. Eine Lösung für dieses Problem wäre das gezielte Sperren von Zeilen im Editor (siehe Abschnitt~\ref{6_2_Perspektiven}).

Für Speicherung und Austausch wählt \verb|javaCodeTransfer.ts| den kürzesten Transfer-String aus drei Verfahren: direkte Wörterbuch-Tokens, deflatierte Tokens oder deflatiertes UTF-8, jeweils gekennzeichnet durch ein Formatpräfix (\texttt{A}, \texttt{B} oder \texttt{C}) mit Base64URL-Nutzdaten. Das integrierte Wörterbuch ersetzt häufige Java-Idiome (z.\ B. \texttt{System.out.println} oder Methodensignaturen) sowie Umlaute durch kompakte 1-Byte-Tokens im Bereich oberhalb des Standard-ASCII-Codes. Falls die Browser-Laufzeit keine Deflate-Kompression unterstützt, stehen die unkomprimierten Tokens als Fallback bereit. Der Export validiert den Quelltext strikt und akzeptiert nur die im Codec zugelassenen Zeichen.
